\chapter{Cofibration categories}
\label{chap:CofibrationAndFibrationCategories}

This chapter defines Anderson-Brown-Cisinski (or ABC) cofibration, fibration and model categories. For simplicity, we refer to ABC cofibration (fibration, model) categories as just cofibration (fibration, model) categories, when no confusion with Quillen model categories or Baues cofibration categories is possible.

What is an ABC cofibration category? It is a category ${\cat M}$ with two distinguished classes of maps, the weak equivalences and the cofibrations, satisfying a set of axioms denoted CF1-CF6. An ABC fibration category is a category ${\cat M}$ with weak equivalences and fibrations, satisfying the dual axioms F1-F6. An ABC model category is a category ${\cat M}$ with weak equivalences, cofibrations and fibrations that is at the same time a cofibration and a fibration category.

Any Quillen model category is an ABC model category (\propositionref{prop:relclosedmodcat}). Diagrams indexed by a small category in an ABC cofibration category form again a cofibration category (\theoremref{thm:generalpointwisecofstructure}), a property not enjoyed in general by Quillen model categories.

The ultimate goal using the cofibration category axioms CF1-CF6 is to construct (in \chapterref{chap:limcolimcofinal}) homotopy colimits in ${\cat M}$ indexed by small categories ${\cat D}$, and more generally to construct 'relative' homotopy colimits along functors $u \colon {\cat D}_1 \ra {\cat D}_2$ of small catgories.

The category ${\cat M}$ will not be assumed in general to be cocomplete. Under the simplifying assumption that ${\cat M}$ is cocomplete, however, the homotopy colimit along a functor $u \colon {\cat D}_1 \ra {\cat D}_2$ is the total left derived functor of $\colim^u \colon {\cat M}^{{\cat D}_1} \ra {\cat M}^{{\cat D}_2}$. It is a known fact that a colimit indexed by a small category ${\cat D}$ can be constructed in terms of pushouts of small sums of objects in ${\cat M}$. Furthermore, a 'relative' colimit $\colim^u$ can be described in terms of absolute colimits indexed by the over categories $(u \downarrow d_2)$ for $d_2 \eps {\cat D}_2$ (see \lemmaref{lem::computerelcolim}).

It is perhaps not surprising then that the cofibration category axioms specify an approximation property of maps by cofibrations (axiom CF4), as well as the behaviour of cofibrations under pushouts (axiom CF3) and under small sums (axiom CF5). The axiom CF6 has a technical rather than conceptual motivation - we simply need it to make the rest of our arguments work (but see \sectionref{sec:oddsandends}).

A large part of the theory can be developed actually from a subset of the axioms, namely the axioms CF1-CF4. The theory of homotopic maps, of homotopy calculus of fractions and of cofibrant approxiomation functors of \chapterref{chap:homotcat} only requires this smaller set of axioms. A category with weak equivalences and cofibrations satisfying the axioms CF1-CF4 will be called a {\it precofibration} category. From the precofibration category axioms, it turns out that one can construct all the homotopy colimits indexed by {\it finite, direct} categories \cite{Cisinski2}. 

One idea worth repeating is that while we need both the cofibrations and the weak equivalences to {\it construct} homotopy colimits, the homotopy colimits are {\it characterized} in the end just by the weak equivalences. So when working with a cofibration category with a fixed class of weak equivalences, it would be desirable to have a class of cofibrations as large as possible. This is where the concept of {\it left proper} maps becomes useful (see \sectionref{sec:properness}). In a {\it left proper} precofibration category $({\cat M}$, ${\cat W}$, ${\cat Cof})$, any cofibration $A \ra B$ with $A$ cofibrant is a left proper map, and the class of left proper maps denoted ${\cat PrCof}$ yields again a precofibration category structure $({\cat M}$, ${\cat W}$, ${\cat PrCof})$. But if ${\cat M}$ is a left proper CF1-CF6 cofibration category, then $({\cat M}$, ${\cat W}$, ${\cat PrCof})$ may not necessarily satisfy the axioms CF5 and CF6. Dual results hold for right proper prefibration categories.

\section{The axioms}
\label{sec:CofibrationCategoriesFibrationCategoriesAxioms}

\index{ABC, or Anderson-Brown-Cisinski!(co)fibration category}
\begin{defn}[Anderson-Brown-Cisinski cofibration categories]
\label{defn:cofcat}
\mbox{} \\
An ABC cofibration category (${\cat M}$, ${\cat W}$, ${\cat Cof}$) consists of a category ${\cat M}$, and two distinguished classes of maps of ${\cat M}$ - the weak equivalences (or trivial maps) ${\cat W}$ and the cofibrations ${\cat Cof}$, subject to the axioms CF1-CF6. The initial object $\initial$ of ${\cat M}$ exists by axiom CF1, and an object $A$ is called {\it cofibrant} if the map $\initial \ra A$ is a cofibration.

The axioms are:
\begin{description}
\item[CF1]
${\cat M}$ has an initial object $\initial$, which is cofibrant. Cofibrations are stable under composition. All isomorphisms of ${\cat M}$ are weak equivalences, and all isomorphisms with the domain a cofibrant object are trivial cofibrations. 

\item[CF2]
\index{axiom!2 out of 3 axiom}(Two out of three axiom) Suppose $f$ and $g$ are maps such that $gf$ is defined. If two of $f$, $g$, $gf$ are weak equivalences, then so is the third.

\item[CF3]
\index{axiom!pushout axiom}(Pushout axiom) Given a solid diagram in ${\cat M}$, with $i$ a cofibration and  $A$, $C$ cofibrant
\begin{center}
$ \xymatrix {
	A \ar[r] \ar@{>->}[d]_i &
	C \ar@{>-->}[d]^j \\
	B \ar@{-->}[r] &
	D
	}$
\end{center}
then
\begin{enumerate}
\item
the pushout exists in ${\cat M}$ and $j$ is a cofibration, and
\item
if additionally $i$ is a trivial cofibration, then so is $j$.
\end{enumerate}

\item[CF4]
\index{axiom!factorization axiom}(Factorization axiom) Any map $f \colon A \ra B$ in ${\cat M}$ with $A$ cofibrant factors as $f = rf'$, with $f'$ a cofibration and $r$  a weak equivalence

\item[CF5]
If $f_i \colon A_i \ra B_i$ for $i \eps I$ is a set of cofibrations with $A_i$ cofibrant, then
\begin{enumerate}
\item
$\Sum A_i$, $\Sum B_i$ exist and are cofibrant, and $\Sum f_i$ is a cofibration.
\item
if additionally all $f_i$ are trivial cofibrations, then so is $\Sum f_i$.
\end{enumerate}

\item[CF6]
For any countable direct sequence of cofibrations with $A_0$ cofibrant
\begin{center}
$ \xymatrix {
        A_0 \ar@{>->}[r]^{a_{0}} &
        A_1 \ar@{>->}[r]^{a_{1}} &
        A_2 \ar@{>->}[r]^{a_{2}} &
	... \\
	}$
\end{center}
\begin{enumerate}
\item
the colimit object $\colim A_n$ exists and the transfinite composition $A_0 \ra \colim A_n$ is a cofibration.
\item
if additionally all $a_i$ are trivial cofibrations, then so is $A_0 \ra \colim A_n$.
\end{enumerate}
\end{description}

If (${\cat M}$, ${\cat W}$, ${\cat Cof}$) only satisfies the axioms CF1-CF4, it is called a precofibration category.
\end{defn}

Pushouts are defined by an universal property, and are only defined up to an unique isomorphism. Since all isomorphisms with cofibrant domain are trivial cofibrations, it does not matter which isomorphic representative of the pushout we choose in CF3.

In the axiom CF4, the map $f^{'}$ is called a {\it cofibrant replacement} of $f$. If $r \colon A^{'} \ra A$ is a weak equivalence with $A^{'}$ cofibrant, the object $A^{'}$ is called a {\it cofibrant replacement} of $A$.

If $A \ra B$ is a cofibration with $A$ cofibrant then $B$ is cofibrant. But there may exist cofibrations $A \ra B$ with $A$ not cofibrant. If we denote ${\cat Cof}^{'}$ the class of cofibrations $A \ra B$ with $A$ cofibrant, then (${\cat M}$, ${\cat W}$, ${\cat Cof}^{'}$) is again a cofibration category. 

We will sometimes refer to a cofibration category as just ${\cat M}$. We will also denote ${\cat M}_{cof}$ \index{${\cat M}_{cof}$, ${\cat M}_{fib}$} the full subcategory of cofibrant objects of ${\cat M}$. 

The category ${\cat M}_{cof}$ is a cofibration category, and in fact so is any full subcategory ${\cat M}^{'}$ of ${\cat M}$ that includes ${\cat M}_{cof}$, with the induced structure (${\cat M}^{'}$, ${\cat W} \cap {\cat M}^{'}$, ${\cat Cof} \cap {\cat M}^{'}$). 

If ${\cat M}$ is a precofibration category, then (${\cat M}$, ${\cat W}$, ${\cat Cof}^{'}$), ${\cat M}_{cof}$ and any ${\cat M}^{'}$ as above are precofibration categories. We will also refer to precofibration categories as CF1-CF4 cofibration categories.

\begin{defn}[Anderson-Brown-Cisinski fibration categories]
\label{defn:fibcat}
\mbox{} \\
An ABC fibration category (${\cat M}$, ${\cat W}$, ${\cat Fib}$) consists of a category ${\cat M}$, and two distinguished classes of maps - the weak equivalences ${\cat W}$ and the fibrations ${\cat Fib}$, subject to the axioms F1-F6. The terminal object $\terminal$ of ${\cat M}$ exists by axiom F1, and an object $A$ of ${\cat Fib}$ is called {\it fibrant} if the map $A \ra \terminal$ is a fibration. 

The axioms are:
\begin{description}
\item[F1]
${\cat M}$ has a final object $\final$, which is fibrant. Fibrations are stable under composition. All isomorphisms of ${\cat M}$ are weak equivalences, and all isomorphisms with fibrant codomain are trivial fibrations. 

\item[F2]
(Two out of three axiom) Suppose $f$ and $g$ are maps such that $gf$ is defined. If two of $f$, $g$, $gf$ are weak equivalences, then so is the third.

\item[F3]
(Pullback axiom) Given a solid diagram in ${\cat M}$, with $p$ a fibration and $A$, $C$ fibrant,
\begin{center}
$ \xymatrix {
	D \ar@{-->}[r] \ar@{-->>}[d]_q &
	B \ar@{->>}[d]^p \\
	C \ar[r] &
	A
	}$
\end{center}
then
\begin{enumerate}
\item
the pullback exists in ${\cat M}$ and $q$ is a fibration, and
\item
if additionally $p$ is a trivial fibration, then so is $q$.
\end{enumerate}

\item[F4]
(Factorization axiom) Any map $f \colon A \ra B$ in ${\cat M}$ with $B$ fibrant factors as $f = f's$, with $s$ a weak equivalence and $f'$  a fibration.

\item[F5]
If $f_i \colon A_i \ra B_i$ for $i \eps I$ is a set of fibrations with $B_i$ fibrant, then 
\begin{enumerate}
\item
$\times A_i$, $\times B_i$ exist and are fibrant, and $\times f_i$ is a fibration
\item
if additionally all $f_i$ are trivial fibrations, then so is $\times f_i$.
\end{enumerate}

\item[F6]
For any countable inverse sequence of fibrations with $A_0$ fibrant
\begin{center}
$ \xymatrix {
    ... \ar@{->>}[r]^{a_{2}} &
    A_2 \ar@{->>}[r]^{a_{1}} &
    A_1 \ar@{->>}[r]^{a_{0}} &
    A_0
  }$
\end{center}
\begin{enumerate}
\item
the limit object $\lim A_i$ exists and the transfinite composition $\lim A_n \ra A_0$ is a fibration
\item
if additionally all $a_i$ are trivial fibrations, then so is $\lim A_n \ra A_0$.
\end{enumerate}
\end{description}

If (${\cat M}$, ${\cat W}$, ${\cat Fib}$) only satisfies the axioms F1-F4, it is called a prefibration category.
\end{defn}

The axioms are dual in the sense that (${\cat M}$, ${\cat W}$, ${\cat Cof}$) is a cofibration category if and only if (${\cat M}^{op}, {\cat W}^{op}, {\cat Cof}^{op}$) is a fibration category. 

In the axiom F4, the map $f^{'}$ is called a {\it fibrant replacement} of $f$. If $r \colon A \ra A^{'}$ is a weak equivalence with $A^{'}$ fibrant, the object $A^{'}$ is called a {\it fibrant replacement} of $A$.

If we denote ${\cat Fib}^{'}$ the class of fibrations $A \ra B$ with $B$ fibrant then (${\cat M}$, ${\cat W}$, ${\cat Fib}^{'}$) again is a fibration category.

We will denote ${\cat M}_{fib}$ to be the full subcategory of fibrant objects of a fibration category ${\cat M}$.  The category ${\cat M}_{fib}$ as well as any full subcategory ${\cat M}^{'}$ of ${\cat M}$ that includes ${\cat M}_{fib}$ satisfy again the axioms of a fibration category. 

If ${\cat M}$ is a prefibration category, then so are (${\cat M}$, ${\cat W}$, ${\cat Fib}^{'}$), ${\cat M}_{fib}$ and any ${\cat M}^{'}$ as above.

\begin{defn}[Anderson-Brown-Cisinski model categories]
\label{defn:abcmodel}
\mbox{} \\
An ABC model category \index{ABC, or Anderson-Brown-Cisinski!model category} $({\cat M}$, ${\cat W}$, ${\cat Cof}, {\cat Fib})$ consists of a category ${\cat M}$ and three distinguished classes of maps ${\cat W}, {\cat Cof}, {\cat Fib}$ with the property that $({\cat M}$, ${\cat W}$, ${\cat Cof})$ is an ABC cofibration category and that $({\cat M}$, ${\cat W}$, ${\cat Fib})$ is an ABC fibration category.
\end{defn}

We will say that ${\cat M}$ is an ABC premodel category if it is only an ABC precofibration and prefibration category.

In all sections of this chapter except \sectionref{sec:transfincomp}, we will do our work assuming only the precofibration category axioms CF1-CF4 (and dually the prefibration category axioms F1-F4). In \sectionref{sec:transfincomp}, we will assume that the full set of axioms is verified.

\section{Sums and products of objects}
\label{sec:sumsproducts}

In general, the objects of a precofibration category are not closed under finite sums. But finite sums of cofibrant objects exist and are cofibrant. Dually, in a prefibration category finite products of fibrant objects exist and are fibrant. In fact we can prove the slightly more general statement:

\begin{lem}
\label{lem:cofsums}
\mbox{}
\begin{enumerate}
\item
Suppose that ${\cat M}$ is a precofibration category. If $f_i \colon A_i \ra B_i$ for $i = 0, ..., n$ are cofibrations with $A_i$ cofibrant, then $\Sum A_i$, $\Sum B_i$ exist and are cofibrant, and $\Sum f_i$ is a cofibration which is trivial if all $f_i$ are trivial.
\item
Suppose that ${\cat M}$ is a prefibration category. If $f_i \colon A_i \ra B_i$ for $i = 0, ..., n$ are fibrations with $B_i$ fibrant, then $\times A_i$, $\times B_i$ exist and are fibrant, and $\times f_i$ is a fibration which is trivial if all $f_i$ are trivial.
\end{enumerate}
\end{lem}

\begin{proof}
We will prove (1), and observe that statement (2) is dual to (1). Using induction on $n$, we can reduce the problem to two maps $f_0 \colon A_0 \ra B_0$ and $f_1 ; A_1 \ra B_1$. If we prove the statement for $f_0, 1_{A_1}$ and $1_{B_0}, f_1$ then the statement follows for $f_0, f_1$. So it suffices to show that if $f \colon A \ra B$ is a (trivial) cofibration and $A$, $C$ are cofibrant, then $A \Sum C$, $B \Sum C$ exist and are cofibrant and $f \Sum 1_C$ is a (trivial) cofibration. From axiom CF3 (1) applied to 
\begin{center}
$ \xymatrix {
	\initial \ar@{>->}[r] \ar@{>->}[d] &
	C \ar@{>->}[d] \\
	A \ar@{>->}[r] &
	A \Sum C
	}$
\end{center}
we see that $A \Sum C$ exists and is cofibrant, and $A \ra A \Sum C$ is a cofibration. Similarly, $B \Sum C$ exists and is cofibrant, and from CF3 applied to
\begin{center}
$ \xymatrix {
	A \ar@{>->}[r] \ar@{>->}[d]_{f} &
	A \Sum C \ar@{>->}[d]^{f \Sum 1_C} \\
	B \ar@{>->}[r] & 
	B \Sum C
	}$
\end{center}
we see that $f \Sum 1_C$ is a cofibration which is trivial if $f$ is trivial.
\end{proof}

\section{Factorization lemmas}
\label{sec:factlemmas}
The Brown Factorization Lemma is an improvement of the factorization axiom CF4 for maps between cofibrant objects.

\index{lemma!Brown factorization}
\begin{lem}[Brown factorization, \cite{Brown}]
\label{lem:brownfact}
\mbox{}
\begin{enumerate}
\item
Let ${\cat M}$ be a precofibration category, and $f \colon A \ra B$ be a map between cofibrant objects. Then $f$ factors as $f = rf'$, where $f'$ is a cofibration and $r$ is a left inverse to a trivial cofibration. 
\item
Let ${\cat M}$ be a prefibration category, and $f \colon A \ra B$ be a map between fibrant objects. Then $f$ factors as $f = f's$, where $f'$ is a fibration and $s$ is a right inverse to a trivial fibration.
\end{enumerate}
\end{lem}

\begin{proof}
The statements are dual, so it suffices to prove (1). We need to construct $f'$, $r$ and $s$ with $f = rf'$ and $rs = 1_B$. 

If we apply the factorization axiom to $f + 1_B$, we get a diagram
\begin{center}
$ \xymatrix {
        A \Sum B \ar[rr]^{f + 1_{B}} \ar@{>->}[rd]_{f^{'} + s} &  & B & \\
        & B^{'} \ar[ru]_{r}^\sim & &
	}$
\end{center}
Since $f^{'} + s$ is a cofibration and $A, B$ are cofibrant, the maps $f^{'}$  and $s$ are cofibrations. The map $r$ is a weak equivalence, and from the commutativity of the diagram we have $rs = 1_B$, therefore $s$ is also a weak equivalence.
\end{proof}

\begin{rem}
\label{rem:brownfact}
We have in fact proved a stronger statement. We have shown that any map $f \colon A \ra B$ between cofibrant objects in a precofibration category factors as $f = rf'$, with $rs = 1_B$ where $f'$, $f^{'} + s$ are cofibrations and $s$ is a trivial cofibration.

Dually, any map $f \colon A \ra B$ between fibrant objects in a prefibration category factors as $f = f's$, with $rs = 1_A$ where $f'$ and $(f^{'}, r)$ are fibrations and $r$ is a trivial fibration.
\end{rem}

Next lemma is a relative version of the factorization axiom CF4 (resp. F4). 

\begin{lem}[Relative factorization of maps]
\label{lem:relativefactorization}
\mbox{}
\begin{enumerate}
\item
Let ${\cat M}$ be a precofibration category, and let
\begin{center}
$\xymatrix {
    A_1 \ar[r]^{f_1} \ar[d]_a & B_1 \ar[d]^b \\
    A_2 \ar[r]^{f_2} & B_2
  }$
\end{center}
be a commutative diagram with $A_1, A_2$ cofibrant. Suppose that $f_1 = r_1 f^{'}_1$ is a factorization of $f_1$ as a cofibration followed by a weak equivalence. Then there exists a commutative diagram
\begin{center}
$\xymatrix{     
    A_1 \ar@{>->}[r]^{f^{'}_1} \ar[d]_a & A^{'}_1 \ar[r]^{r_1}_\sim \ar[d]_(.44){a^{'}} & B_1 \ar[d]^b \\
    A_2 \ar@{>->}[r]^{f^{'}_2} & A^{'}_2 \ar[r]^{r_2}_\sim & B_2
}$ 
\end{center}
where $r_2 f^{'}_2$ is a factorization of $f_2$ as a cofibration followed by a weak equivalence and such that $A_2 \Sum_{A_1} A^{'}_1 \ra A^{'}_2$ is a cofibration.
\item
%Let ${\cat M}$ be a prefibration category, and let
%\begin{center}
%$\xymatrix {
%    A_1 \ar[r]^{f_1} \ar[d]_a & B_1 \ar[d]^b \\
%    A_2 \ar[r]^{f_2} & B_2
%  }$
%\end{center}
%be a commutative diagram with $B_1, B_2$ fibrant. Suppose that $f_2 = f^{'}_2 s_2$ is a factorization of $f_2$ as a weak equivalence followed by a fibration. Then there exists a commutative diagram
%\begin{center}
%$\xymatrix{     
%    A_1 \ar[r]^{s_1}_\sim \ar[d]_a & B^{'}_1 \ar@{->>}[r]^{f^{'}_1} \ar[d]_(.44){b^{'}} & B_1 \ar[d]^b \\
%    A_2 \ar[r]^{s_2}_\sim & B^{'}_2 \ar@{->>}[r]^{f^{'}_2} & B_2
%}$ 
%\end{center}
%where $f^{'}_1 s_1$ is a factorization of $f_1$ as a weak equivalence followed by a fibration and such that $B^{'}_1 \ra B^{'}_2 \times_{B_2} B_1$ is a fibration.
The dual of (1) holds for prefibration categories.
\end{enumerate}
\end{lem}

\begin{proof}
To prove (1), in the commutative diagram
\begin{center}
$\xymatrix {
    A_1 \ar[d]_{a} \ar@{>->}[r]^{f^{'}_1} &
    A^{'}_1 \ar[d] \ar[rr]_{\sim}^{r_1} & &
    B_1 \ar[d]^b \\
    A_2 \ar@{>->}[r] &
    A_2 \Sum_{A_1} A^{'}_1 \ar@{>->}[r]^-s &
    A^{'}_2 \ar[r]_{\sim}^{r_2} &
    B_2
  }$
\end{center}
the pushout $A_2 \Sum_{A_1} A^{'}_1$ exists by CF3, and we construct the cofibration $s$ and the weak equivalence $r_2$ using the factorization axiom CF4 applied to $A_2 \Sum_{A_1} A^{'}_1 \lra B_2$. 
\end{proof}

The Brown Factorization Lemma has the following relative version:

\begin{lem}[Relative Brown factorization]
\label{lem:relbrownfact}
\mbox{}
\begin{enumerate}
\item
Suppose that ${\cat M}$ is a precofibration category, and that
\begin{center}
$ \xymatrix {
        A_1 \ar[r]^{f_1} \ar[d]_a &
	B_1 \ar[d]^b \\
	A_2 \ar[r]^{f_2} &
	B_2
	}$
\end{center}
is a commutative diagram with cofibrant objects. Suppose that $f_1 = r_1f^{'}_1$, $r_1 s_1 = 1$ is a Brown factorization of $f_1$, with $f^{'}_1$, $f^{'}_1 + s_1$ cofibrations and $s_1$ a trivial cofibration. Then there exists a Brown factorization $f_2 = r_2f^{'}_2$, $r_2 s_2 = 1 $ with $f^{'}_2$, $f^{'}_2 + s_2$ cofibrations and $s_2$ a trivial cofibration and a map $b^{'}$ such that in the diagram
\begin{center}
$ \xymatrix {
        A_1 \ar@{>->}[r]^{f^{'}_1} \ar[d]_a &
        B^{'}_1 \ar@<1ex>[r]^{r_1} \ar[d]^{b^{'}} &
	B_1 \ar[d]^b \ar@<1ex>@{>->}[l]^{s_1}_{\sim} \\
	A_2 \ar@{>->}[r]^{f^{'}_2} &
	B^{'}_2 \ar@<1ex>[r]^{r_2} &
	B_2 \ar@<1ex>@{>->}[l]^{s_2}_{\sim}
	}$
\end{center}
we have that $b^{'}f^{'}_1 = f^{'}_2 a$, $br_1 = r_2 b^{'}$, $b^{'}s_1 = s_2 b$,  and that
\begin{center}
$A_2 \Sum_{A_1} B^{'}_1 \ra B^{'}_2$ and $B_2 \Sum_{B_1} B^{'}_1 \ra B^{'}_2$
\end{center}
are a cofibration (resp. a trivial cofibration).
\item
The dual of (1) holds for prefibration categories.
\end{enumerate}
\end{lem}

\begin{proof}
To prove (1), denote $A_3 = A_2 \Sum_{A_1} B^{'}_1$ and $B_3 = B_2 \Sum_{B_1} B^{'}_1$. 
\begin{center}
$ \xymatrix {
        A_1 \ar@{>->}[rr]^{f^{'}_1} \ar[d]_a &&
        B^{'}_1 \ar@<1ex>[rr]^{r_1} \ar[dr] \ar[dl] &&
	B_1 \ar[d]^b \ar@<1ex>@{>->}[ll]^{s_1}_{\sim} \\
	A_2 \ar@{>->}[r]^{f^{'}} &
	A_3 \ar[rr]^{f_3} && B_3 \ar@<1ex>[r]^{r} &
	B_2 \ar@<1ex>@{>->}[l]^{s}_{\sim}
	}$
\end{center}
We apply \lemmaref{lem:relativefactorization} to the commutative diagram
\begin{center}
$\xymatrix {
    A_1 \Sum B_1 \ar@{>->}[r]^-{f^{'}_1 + s_1} \ar[d]_{f^{'} a \Sum s b} & B^{'}_1 \ar[r]^{r_1}_\sim & B_1 \ar[d]^{sb} \\
    A_3 \Sum B_3 \ar[rr]^-{f_3 + 1_{B_3}} & & B_3
  }$
\end{center}
and we construct a commutative diagram
\begin{center}
$ \xymatrix {
        A_1 \Sum B_1 \ar[rr]^{f_1 + 1_{B_1}} \ar[dd]_{f^{'} a \Sum s b} \ar@{>->}[rd]_{f^{'}_1 + s_1} &  & B_1 \ar[dd]^{sb} & \\
        & B^{'}_1 \ar[ru]_{r_1}^\sim \ar[dd]^(.3){b^{'}} & & \\
        A_3 \Sum B_3 \ar'[r][rr] \ar@{>->}[rd]_{f^{'}_3 + s_3} & & B_3 & \\
        & B^{'}_2 \ar[ru]_{r_3}^\sim & &
	}$
\end{center}
We now set $f^{'}_2 = f^{'}_3f^{'}$, $r_2 = rr_3$ and $s_2 = s_3s$.
\end{proof}

\section{Extension lemmas}
\label{sec:extlemmas}
The Gluing Lemma describes the behavior of cofibrations and weak equivalences under pushouts, and is one of the basic building blocks we will employ in the construction of homotopy colimits.

\index{lemma!gluing lemma}
\begin{lem}[Gluing Lemma]
\label{lem:gluinglemma1}
\mbox{}
\begin{enumerate}
\item
Let ${\cat M}$ be a precofibration category. In the diagram
\begin{center}
$ \xymatrix {
        A_1 \ar@{>->}[rr]^{f_{12}} \ar[dd]_{u_1} \ar[rd]_{f_{13}} &  & A_2 \ar'[d][dd]^{u_2} \ar[rd] & \\
        & A_3 \ar@{>->}[rr] \ar[dd]^(.3){u_3} & & A_4 \ar[dd]^{u_4} \\
        B_1 \ar@{>->}'[r][rr]_{g_{12}} \ar[rd]_{g_{13}} & & B_2 \ar[rd] & \\
        & B_3 \ar@{>->}[rr] & & B_4
	}$
\end{center}
suppose that $A_1$, $A_3$, $B_1$, $B_3$ are cofibrant, that $f_{12}, g_{12}$ are cofibrations, and that the top and bottom faces are pushouts. 
\begin{enumerate}
\item
If $u_1, u_3$ are cofibrations and the natural map \\ $B_1 \Sum_{A_1} A_2 \lra B_2$ is a cofibration, then $u_2, u_4$ and the natural map $B_3 \Sum_{A_3} A_4 \lra B_4$ are cofibrations.
\item
If $u_1, u_2, u_3$ are weak equivalences, then $u_4$ is a weak equivalence.
\end{enumerate}
\item
Let ${\cat M}$ be a prefibration category. In the diagram
\begin{center}
$ \xymatrix {
        B_4 \ar@{->>}[rr] \ar[dd]_{u_4} \ar[rd] &  & B_3 \ar'[d][dd]^{u_3} \ar[rd]^{g_{31}} & \\
        & B_2 \ar@{->>}[rr]^(.3){g_{21}} \ar[dd]^(.3){u_2} & & B_1 \ar[dd]^{u_1} \\
        A_4 \ar@{->>}'[r][rr] \ar[rd] & & A_3 \ar[rd]^{f_{31}} & \\
        & A_2 \ar@{->>}_{f_{21}}[rr] & & A_1
	}$
\end{center}
suppose that $A_1$, $A_3$, $B_1$, $B_3$ are fibrant, that $f_{21}, g_{21}$ are fibrations, and that the top and bottom faces are pullbacks. 
\begin{enumerate}
\item
If $u_1, u_3$ are fibrations and the natural map \\ $B_2 \lra B_1 \times_{A_1} A_2 $ is a fibration, then $u_2, u_4$ and the natural map $B_4 \lra B_3 \times_{A_3} A_4$ are fibrations.
\item
If $u_1, u_2, u_3$ are weak equivalences, then $u_4$ is a weak equivalence.
\end{enumerate}
\end{enumerate}
\end{lem}

\begin{proof}
The statements are dual, and we will prove only (1). Let $B^{'}_2 = B_1 \Sum_{A_1} A_2$ and $B^{'}_4 = B_3 \Sum_{A_3} A_4$ be the pushout of the front and back faces of the diagram of (1). These pushouts exist because of the pushout axiom CF3.
\begin{center}
$ \xymatrix {
        A_1 \ar@{>->}[rr]^{f_{12}} \ar[ddd]_(.4){u_1} \ar[rd]_{f_{13}} & & A_2 \ar[dd]^(.7){u^{'}_2} \ar[rd] & \\
        & A_3 \ar@{>->}[rr] \ar[ddd]^(.2){u_3} & & A_4 \ar[dd]^(.35){u^{'}_4} \\
        & & B^{'}_2 \ar[d]^{u^{''}_2} \ar[rd] & \\
        B_1 \ar@{>->}[rr]_(.7){g_{12}} \ar[rd]_{g_{13}} \ar@{>->}[rru]^(.3){g^{'}_{12}} & & B_2 \ar[rd] & B^{'}_4 \ar[d]^{u^{''}_4} \\
        & B_3 \ar@{>->}[rr] \ar@{>->}[rru]^(.3){g^{'}_{34}} & & B_4
	}$
\end{center}
The maps $g^{'}_{12}$ and $g^{'}_{34}$ are pushouts of cofibrations, therefore cofibrations. Furthermore, we observe that $B^{'}_4 = B_3 \Sum_{B_1} B^{'}_2$, and therefore $u^{''}_4$ is the pushout of $u^{''}_2$ along $B^{'}_2 \ra B^{'}_4$. 

Let's prove (a). If $u_1, u_3$ and $u^{''}_2$ are cofibrations, then by the pushout axiom $u^{'}_2$ and $u^{'}_4$ are cofibrations. $u^{''}_4$ is a pushout of $u^{''}_2$, therefore also a cofibration. It follows that $u_2, u_4$ are cofibrations.

Let's now prove (b), first under the assumption that
\begin{equation}
\label{eqn:gluinglemmasimplifying}
u_1, u_3 \;\; \mathrm{ and } \;\; u^{''}_2 \;\; \mathrm{ are} \;\; \mathrm{ cofibrations }
\end{equation}
If they are, since $u_1$, $u_2$, $u_3$ are weak equivalences we see that $u_1$, $u_3$ and their pushouts $u^{'}_2$, $u^{'}_4$ must be trivial cofibrations. From the two out of three axiom, $u^{''}_2$ is a weak equivalence, therefore a trivial cofibration, and so its pushout $u^{''}_4$ also is a trivial cofibration, which shows that $u_4 = u^{''}_4u^{'}_4$ is a weak equivalence.

For general weak equivalences $u_1$, $u_2$, $u_3$, we use the relative Brown factorization lemma to construct the diagram
%\begin{center}
%$ \xymatrix {
%        A_1 \ar@{>->}[rr]^-{f_{12}} \ar@{>->}[dd]_{v_1} \ar[rd]_{f_{13}} &  & A_2 \ar@{>->}[dd]^{v_2} & \\
%        & A_3 \ar@{>->}[dd]^(.3){v_3} & & \\
%        B^{'}_1 \ar@{>->}'[r]^{h_{12}}[rr] \ar@<1ex>[dd]^{r_1}_\sim \ar[rd]_{h_{13}} &  & B^{'}_2 \ar@<1ex>[dd]^{r_2}_\sim & \\
%        & B^{'}_3 \ar@<1ex>[dd]^(.3){r_3}_(.3)\sim & & \\
%        B_1 \ar@{>->}'[r][rr]_{g_{12}} \ar[rd]_{g_{13}} \ar@{>->}@<1.3ex>[uu]^{w_1} & & B_2 \ar@{>->}@<1.3ex>[uu]^{w_2} & \\
%        & B_3 \ar@{>->}@<1.3ex>[uu]^(.7){w_3} & &
%	}$
%\end{center}
\begin{center}
$ \xymatrix {
    A_2 \ar@{>->}[r]^{v_2} & B^{'}_2 \ar@<1ex>[r]^{r_2}_\sim & B_2 \ar@{>->}@<1ex>[l]^{w_2} \\
    A_1 \ar@{>->}[u]^{f_{12}} \ar[d]_{f_{13}} \ar@{>->}[r]^{v_1} & 
    B^{'}_1 \ar@{>->}[u]^{h_{12}} \ar[d]_{h_{13}} \ar@<1ex>[r]^{r_1}_\sim & 
    B_1 \ar@{>->}[u]_{g_{12}} \ar[d]^{g_{13}} \ar@{>->}@<1ex>[l]^{w_1} \\
    A_3 \ar@{>->}[r]_{v_3} & B^{'}_3 \ar@<1ex>[r]^{r_3}_\sim & B_3 \ar@{>->}@<1ex>[l]^{w_3} \\
	}$
\end{center}
In this diagram:
\begin{enumerate}
\item
$v_1$, $w_1$, $r_1$ are constructed as a Brown factorization of $u_1$ as in \remarkref{rem:brownfact}
\item
$v_i$, $w_i$, $r_i$ for $i = 2, 3$ are constructed as relative Brown factorizations of $u_2$ resp. $u_3$ over the Brown factorization $v_1$, $w_1$, $r_1$.
\end{enumerate}

The maps ($w_1$, $w_2$, $w_3$) are trivial cofibrations and ($u_1$, $u_2$, $u_3$) are weak equivalences, so ($v_1$, $v_2$, $v_3$) are trivial cofibrations.

Statement (b) is true for 
\begin{enumerate}
\item[-]
($v_1$, $v_2$, $v_3$) resp. ($w_1$, $w_2$, $w_3$) because they satisfy property (\ref{eqn:gluinglemmasimplifying})
\item[-]
therefore true for ($r_1$, $r_2$, $r_3$) as a left inverse to ($w_1, w_2, w_3$)
\item[-]
therefore true for ($u_1$, $u_2$, $u_3$) as the composition of ($v_1$, $v_2$, $v_3$) and ($r_1$, $r_2$, $r_3$).
\end{enumerate}
\end{proof}

As a corollary we have

\index{excision}
\index{lemma!excision lemma}
\begin{lem}[Excision]
\label{lem:excision}
\mbox{}
\begin{enumerate}
\item
Let ${\cat M}$ be a precofibration category. In the diagram below
\begin{center}
$ \xymatrix {
        A \ar[r]^f_{\sim} \ar@{>->}[d]_i & C \\
	B & 
	}$
\end{center}
suppose that $A$, $C$ are cofibrant, $i$ is a cofibration and $f$ is a weak equivalence. Then the pushout of $f$ along $i$ is again a weak equivalence.
\item
Let ${\cat M}$ be a prefibration category. In the diagram below
\begin{center}
$ \xymatrix {
        & B \ar@{->>}[d]^p \\
        C \ar[r]^f_{\sim} & A
	}$
\end{center}
suppose that $A$, $C$ are fibrant, $p$ is a fibration and $f$ is a weak equivalence. Then the pullback of $f$ along $p$ is again a weak equivalence.
\end{enumerate}
\end{lem}

\begin{proof}
Part (1) is a particular case of the Gluing Lemma (1) for $f_{12} = g_{12} = i$, $u_1 = f_{13} = 1_A$, $g_{13} = u_3 = f$ and $u_2 = 1_B$. Part (2) is dual.
\end{proof}

It is now easy to see that in the presence of the rest of the axioms, the axiom CF3 (2) is equivalent to the Gluing Lemma (1) (b) and to excision. A dual statement holds for the fibration axioms.

\begin{lem}[Equivalent formulation of CF3]
\label{lem:equivaxiomcf3}
\mbox{}
\begin{enumerate}
\item
If $({\cat M}$, ${\cat W}$, ${\cat Cof})$ satisfies the axioms CF1-CF2, CF3 (1) and CF4, then the following are equivalent:
\begin{enumerate}
\item
It satisfies CF3 (2)
\item
It satisfies the Gluing Lemma (1) (b)
\item
It satisfies the Excision \lemmaref{lem:excision} (1).
\end{enumerate}
\item
If $({\cat M}$, ${\cat W}$, ${\cat Fib})$ satisfies the axioms F1-F2, F3 (1) and F4, then the following are equivalent:
\begin{enumerate}
\item
It satisfies F3 (2)
\item
It satisfies the Gluing Lemma (2) (b)
\item
It satisfies the Excision \lemmaref{lem:excision} (2).
\end{enumerate}
\end{enumerate}
\end{lem}

\begin{proof}
It suffices to prove (1). We have proved (a) $\Ra$ (b) as \lemmaref{lem:gluinglemma1}, and we have seen (b) $\Ra$ (c) in the proof of \lemmaref{lem:excision}.

For (c) $\Ra$ (a), suppose that $A, C$ are cofibrant, that $i \colon A \rightarrowtail B$ is a trivial cofibration and that $f \colon A \ra C$ is a map. Factor $f$ as a cofibration $f^{'}$ followed by a weak equivalence $r$, and using axiom CF3 (1) construct the pushouts $g^{'}, s$ of $f^{'}, r$
\begin{center}
$ \xymatrix {
    A \ar@{>->}[r]^{f^{'}} \ar@{>->}[d]_{i}^{\sim} & C^{'} \ar[r]^r_\sim \ar@{>->}[d]^{i^{'}}_\sim
 & C \ar@{>->}[d]^j \\
    B \ar@{>->}[r]^{g^{'}} & B^{'} \ar[r]^{s}_\sim & D
  }$
\end{center}
The maps $g^{'}, i^{'}$ and $j$ are cofibrations by axiom CF3 (1). Using excision, $i^{'}$ and $s$ are weak equivalences, so by the 2 out of 3 axiom CF2 the map $j$ is also a weak equivalence.
\end{proof}

\section{Cylinder and path objects}
\label{sec:cylobj}
We next define cylinder objects in a precofibration category, and show that cylinder objects exist. Dually, we define and prove existence of path objects in a prefibration category. 

\index{cylinder object}
\index{path object}
\begin{defn}[Cylinder and path objects]
\label{defn:cylobj}
\mbox{}
\begin{enumerate}
\item
Let ${\cat M}$ be a precofibration category, and $A$ a cofibrant object of ${\cat M}$. A cylinder object for $A$ consists of an object $IA$ and a factorization of the codiagonal $ \nabla \colon \xymatrix { A \Sum A \ar@{>->}[r]^-{i_0+i_1} & IA \ar[r]^{p}_{\sim} & A } $, with $i_0+i_1$ a cofibration and $p$ a weak equivalence.

\item
Let ${\cat M}$ be a prefibration category, and $A$ a fibrant object of ${\cat M}$. A path object for $A$ consists of an object $A^I$ and a factorization of the diagonal $ \Delta \colon \xymatrix { A \ar[r]^{i}_{\sim} & A^I \ar@{->>}[r]^-{(p_0,p_1)} & A \times A } $, with $(p_0,p_1)$ a fibration and $i$ a weak equivalence.
\end{enumerate}
\end{defn}

\begin{lem}[Existence of cylinder and path objects]
\label{lem:existencecyl}
\mbox{}
\begin{enumerate}
\item
Let ${\cat M}$ be a precofibration category, and $A$ a cofibrant object of ${\cat M}$. Then $A$ admits a (non-functorial) cylinder object.

\item
Let ${\cat M}$ be a prefibration category, and $A$ a fibrant object of ${\cat M}$. Then $A$ admits a (non-functorial) path object.
\end{enumerate}
\end{lem}

\begin{proof}
To prove (1), observe that if $A$ is cofibrant then the sum $A \Sum A$ exists and is cofibrant by \lemmaref{lem:cofsums}, and we can then use the factorization axiom CF4 to construct a cylinder object $ \xymatrix { A \Sum A \ar@{>->}[r]^-{i_0+i_1} & IA \ar[r]^{p}_{\sim} & A } $. The statement (2) follows from duality.
\end{proof}

Observe that for cylinder objects, the inclusion maps $i_0,i_1 \colon A \lra IA$ are trivial cofibrations. For path objects, the projection maps $p_0,p_1 \colon A^I \lra A$ are trivial fibrations.

\begin{lem}[Relative cylinder and path objects]
\label{lem:existencerelcyl}
\mbox{}
\begin{enumerate}
\item
Let ${\cat M}$ be a precofibration category, and $f \colon A \lra B$ a map with $A, B$ cofibrant objects. Let $IA$ be a cylinder of $A$. Then there exists a cylinder $IB$ and a commutative diagram
\begin{center}
$\xymatrix {
    A \Sum A \ar[d]_{f \Sum f} \ar@{>->}[r] &
    IA \ar[d]_{If} \ar[r]^{\sim} &
    A \ar[d]^f \\
    B \Sum B \ar@{>->}[r] &
    IB \ar[r]^{\sim} &
    B
  }$
\end{center}
with $ \xymatrix { (B \Sum B) \Sum_{A \Sum A} IA \ar@{>->}[r] & IB } $ a cofibration.
\item
Let ${\cat M}$ be a prefibration category,  and $f \colon A \lra B$ a map with $A, B$ fibrant objects. Let $B^I$ be a path object for of $B$. Then there exists a path object $A^I$ and a commutative diagram
\begin{center}
$ \xymatrix {
	A \ar[d]_f \ar[r]^{\sim} & 
	A^I \ar[d]_{f^I} \ar@{->>}[r] &
        A \times A \ar[d]^{f \times f} \\
	B \ar[r]^{\sim} &
        B^I \ar@{->>}[r] &
        B \times B
	}$
\end{center}
with $ \xymatrix { A^I  \ar@{->>}[r] & B^I \times_{B \times B} (A \times A) } $ a fibration.
\end{enumerate}
\end{lem}

\begin{proof}
To prove (1), apply \lemmaref{lem:relativefactorization} to the diagram 
\begin{center}
$\xymatrix {
    A \Sum A \ar[d]_{f \Sum f} \ar@{>->}[r] &
    IA  \ar[r]^{\sim} &
    A \ar[d]^f \\
    B \Sum B \ar[rr]^{\nabla} && B
  }$
\end{center}
Statement (2) is dual to (1).
\end{proof}


\section{Elementary consequences of CF5 and CF6}
\label{sec:transfincomp}
In the previous sections we have proved a number of elementary lemmas that are consequences of the precofibration category axioms CF1-CF4. In this section, we will do the same bringing in one by one the cofibration category axioms CF5 and CF6.

A word on the motivation behind the two additional axioms CF5-CF6. In the construction of homotopy colimits indexed by small diagrams in a cofibration category, it turns out that the role of small {\it direct categories} (\definitionref{def:directinvcat}) is essential, because an arbitrary small diagram can be approximated by a diagram indexed by a small direct category (\sectionref{sec:colimarbitrarycat}). 

For a small direct category ${\cat D}$, its degreewise filtration can be used to show that colimits indexed by ${\cat D}$ may be constructed using small sums, pushouts and countable direct transfinite compositions (at least if the base category is cocomplete). To put things in perspective, the axiom CF3 is a property of pushouts, the axiom CF5 is a property of small sums of maps and the axiom CF6 is a property of countable direct transfinite compositions of maps.

Let us clarify for a moment what we mean in CF6 and F6 by transfinite direct and inverse compositions of maps.
\begin{defn}
\label{defn:transcomp}
Let ${\cat M}$ be a category, and let $k$ be an ordinal. 
\begin{enumerate}
\item
A {\it direct $k$-sequence} of maps (or a {\it direct sequence of length $k$})
\begin{center}
$ \xymatrix {
        A_0 \ar[r]^{a_{01}} &
        A_1 \ar[r]^{a_{12}} &
	... \ar[r] &
        A_i \ar[r] &
        ... &
        (i < k)
	}$
\end{center}
consists of a collection of objects $A_i$ for $i < k$ and maps $a_{i_1i_2} \colon A_{i_1} \ra A_{i_2}$ for $i_1 < i_2 < k$, such that $a_{i_2i_3}a_{i_1i_2} = a_{i_1i_3}$ for all $i_1 < i_2 < i_3 < k$. The map $A_0 \ra \colim^{i < k} A_i$, if the colimit exists, is called the {\it transfinite composition} of the direct $k$-sequence. 
\item
An {\it inverse $k$-sequence} of maps (or an {\it inverse sequence of maps of length $k$})
\begin{center}
$ \xymatrix {
    ... \ar[r] &
    A_i \ar[r] &
    ... \ar[r]^{a_{21}} &
    A_1 \ar[r]^{a_{10}} &
    A_0 &
    (i < l)
  }$
\end{center}
consists of a collection of objects $A_i$  for $i < k$ and maps $a_{i_2i_1} \colon A_{i_2} \ra A_{i_1}$ for $i_1 < i_2 < k$, such that $a_{i_2i_1}a_{i_3i_2} = a_{i_3i_1}$ for all $i_1 < i_2 < i_3 < k$. The map $\lim^{i < k} A_i \ra A_0$, if the limit exists, is called the transfinite composition of the inverse $k$-sequence.
\end{enumerate}
\end{defn}

A direct $k$-sequence of maps is nothing but a diagram in ${\cat M}$ indexed by ordinals $< k$. A map of direct $k$-sequences is a map of such diagrams. 
 
If ${\cat M}$ is a precofibration category, a direct $k$-sequence of (trivial) cofibrations is a direct $k$-sequence in which all maps $a_{ij}$, $i < j < k$ are (trivial) cofibrations.  If ${\cat M}$ is a fibration category, an inverse $k$-sequence of (trivial) fibrations is an inverse $k$-sequence in which all maps $a_{ji}$, $i < j < k$ are (trivial) fibrations.

A sufficient condition for the additional axioms CF5-CF6 to be satisfied is that cofibrations and trivial cofibrations with cofibrant domain are stable under all {\it small} transfinite direct compositions.

\begin{lem}
\label{lem:transfinitecompositionstability}
\mbox{}
\begin{enumerate}
\item
If $({\cat M}, {\cat W}, {\cat Cof})$ satisfies axioms CF1-CF4, and if cofibrations (resp. trivial cofibrations) with cofibrant domain are stable under transfinite compositions of direct $k$-sequences for any small ordinal $k$, then $({\cat M}, {\cat W}, {\cat Cof})$ also satisfies axioms CF5-CF6.
\item
If $({\cat M}, {\cat W}, {\cat Fib})$ satisfies axioms F1-F4, and if fibrations (resp. trivial fibrations) with fibrant codomain are stable under transfinite compositions of inverse $k$-sequences for any small ordinal $k$, then $({\cat M}, {\cat W}, {\cat Fib})$ also satisfies axioms F5-F6.
\end{enumerate}
\end{lem}

\begin{proof}
We only prove (1). Axiom CF6 is clearly verified for ${\cat M}$, since it states that cofibrations (resp. trivial cofibrations) with cofibrant domain are stable under {\it countable} transfinite direct compositions. For the axiom CF5, given a set of (trivial) cofibrations $f_i \colon A_i \ra B_i$ for $i \eps I$ with $A_i$ cofibrant, we choose a well ordering of $I$. Denote $I^+$ the well ordered set $I$ with a maximal element adjoined. $I^+$ can be viewed as the succesor ordinal of $I$, and all elements $i \eps I$ can be viewed as the ordinals smaller than $I$.

We show that for any $i \eps I^+$, we have that $\Sum_{_{k < i}} f_k \colon \Sum_{_{k < i}} A_k \ra \Sum_{_{k < i}} B_k$ is well defined and is a (trivial) cofibration with a cofibrant domain. We use transfinite induction, and the initial step is trivial. Suppose the statement is true for all elements $<i$, and let's prove it for $i$.

If $i$ is a succesor ordinal, the statement for $i$ follows from \lemmaref{lem:cofsums}. Suppose that $i$ is a limit ordinal. 

For any $i^{''} < i^{'} < i$, the inclusion
\[
\Sum_{_{k < i^{''}}} A_k \ra \Sum_{_{k < i^{'}}} A_k
\]
is a cofibration, using the inductive hypothesis. The transfinite composition of these cofibrations defines $\Sum_{_{k < i}} A_k$, which therefore exists and is cofibrant. Similarly, $\Sum_{_{k < i}} B_k$ exists and is cofibrant.

For any $i^{''} < i^{'} < i$, the map
\[
\Sum_{_{k < i^{''}}} B_k \Sum_{_{i^{''} \le k < i}} A_k \ra \Sum_{_{k < i^{'}}} B_k \Sum_{_{i^{'} \le k < i}} A_k
\]
given by
\[
(\Sum_{_{k < i^{''}}} 1_{B_k}) \Sum (\Sum_{_{i^{''} \le k < i^{'}}} f_k) \Sum (\Sum_{_{i^{'} \le k < i}} 1_{A_k}) 
\]
is a well defined (trivial) cofibration with cofibrant domain, using the inductive hypothesis. The transfinite composition of these (trivial) cofibrations with cofibrant domain defines $\Sum_{_{k < i}} f_k$, which is therefore a (trivial) cofibration. The statement of our lemma now follows if we take $i$ to be the maximal element of $I^+$.
\end{proof}

The next two lemmas describe properties of the additional axiom CF5 (resp. F5).

\begin{lem}
\label{lem:smallcofsums}
\mbox{}
\begin{enumerate}
\item
Suppose that ${\cat M}$ is a CF1-CF4 (resp. CF1-CF5) cofibration category. If $f_i \colon A_i \ra B_i$ for $i \eps I$ is a finite (resp. small) set of weak equivalences between cofibrant objects, then $\Sum f_i$ is a weak equivalence.
\item
Suppose that ${\cat M}$ is an F1-F4 (resp. F1-F5) fibration category. If $f_i \colon A_i \ra B_i$ for $i \eps I$ is a finite (resp. small) set of weak equivalences between fibrant objects, then $\times f_i$ is a weak equivalence.
\end{enumerate}
\end{lem}

\begin{proof}
We only prove (1). Using the Brown Factorization Lemma, write $f_i = r_if^{'}_i $, where $f^{'}_i$ is a trivial cofibration and $r_i$ is a left inverse to a trivial cofibration $s_i$. Under both alternative hypotheses, the maps $\Sum f^{'}_i$ and $\Sum s_i$ are trivial cofibrations, so $\Sum r_i$ and therefore $\Sum f_i$ are weak equivalences.
\end{proof}

\begin{lem}[Equivalent formulation of CF5]
\label{lem:equivformaxiomcf5}
\mbox{}
\begin{enumerate}
\item
Suppose that $({\cat M}$, ${\cat W}$, ${\cat Cof})$ satisfies axioms CF1-CF4 and CF5 (1). Then the following are equivalent:
\begin{enumerate}
\item
It satisfies axiom CF5 (2)
\item
The class of weak equivalences between cofibrant objects is stable under small sums.
\end{enumerate}
\item
Suppose that $({\cat M}$, ${\cat W}$, ${\cat Fib})$ satisfies axioms F1-F4 and F5 (1). Then the following are equivalent:
\begin{enumerate}
\item
It satisfies axiom F5 (2)
\item
The class of weak equivalences between fibrant objects is stable under small products.
\end{enumerate}
\end{enumerate}
\end{lem}

\begin{proof}
The implication $(a) \Ra (b)$ is a consequence of \lemmaref{lem:smallcofsums}, and $(b) \Ra (a)$ is trivial.
\end{proof}

\comment{
We have seen from the CF1-CF4 axioms that (trivial) cofibrations with cofibrant domain are stable under {\it finite} sums (this is \lemmaref{lem:cofsums}). We can improve on that using the axioms CF1-CF5 as follows.

\begin{lem}
\label{lem:partialCF6}
\mbox{}
\begin{enumerate}
\item
Let ${\cat M}$ be a CF1-CF5 cofibration category. Then (trivial) cofibrations with cofibrant domain are stable under {\it countable} sums.
\item
Let ${\cat M}$ be an F1-F5 fibration category. Then (trivial) fibrations with cofibrant domain are stable under {\it countable} products.
\end{enumerate}
\end{lem}

\begin{proof}
The proof is very similar to that of \lemmaref{lem:transfinitecompositionstability}, and is omitted.
\end{proof}
} %comment

Here is a consequence of the axiom CF6 (resp. F6).

\begin{lem}
\label{lem:transcompequiv}
\mbox{}
\begin{enumerate}
\item
Let ${\cat M}$ be a CF1-CF4 cofibration category satisfying CF6. For any map of countable direct sequences of cofibrations with $A_0$, $B_0$ cofibrant and all $f_n$ weak equivalences
\begin{center}
$ \xymatrix {
        A_0 \ar@{>->}[r]^{a_{0}} \ar[d]_{f_0}^{\sim} &
        A_1 \ar@{>->}[r]^{a_{1}} \ar[d]_{f_1}^{\sim} &
        A_2 \ar@{>->}[r]^{a_{2}} \ar[d]_{f_2}^{\sim} &
	... \\
        B_0 \ar@{>->}[r]^{b_{0}} &
        B_1 \ar@{>->}[r]^{b_{1}} &
        B_2 \ar@{>->}[r]^{b_{2}} &
	...
	}$
\end{center}
the colimit map $\colim f_n \colon \colim A_n \ra \colim B_n$ is a weak equivalence between cofibrant objects.
\item
Let ${\cat M}$ be an F1-F4 fibration category satisfying CF6. For any map of countable inverse sequences of fibrations with $A_0$, $B_0$ fibrant and all $f_n$ weak equivalences
\begin{center}
$ \xymatrix {
    ... \ar@{->>}[r]^{a_{2}} & \
    A_2 \ar@{->>}[r]^{a_{1}} \ar[d]^{f_2}_{\sim} &
    A_1 \ar@{->>}[r]^{a_{0}} \ar[d]^{f_1}_{\sim} &
    A_0 \ar[d]^{f_0}_{\sim} \\
    ... \ar@{->>}[r]^{b_{2}} &
    B_2 \ar@{->>}[r]^{b_{1}} &
    B_1 \ar@{->>}[r]^{b_{0}} &
    B_0
  }$
\end{center}
the limit map $\lim f_n \colon \lim A_n \ra \lim B_n$ is a weak equivalence between fibrant objects.
\end{enumerate}
\end{lem}

\begin{proof}
To prove (1), observe that using Rem. \ref{rem:brownfact} and \lemmaref{lem:relbrownfact}, we can inductively construct Brown factorizations of $f_n$ that make the diagram
\begin{center}
$ \xymatrix {
        A_0 \ar@{>->}[r]^{a_{0}} \ar@{>->}[d]_{f_0^{'}}^{\sim} &
        A_1 \ar@{>->}[r]^{a_{1}} \ar@{>->}[d]_{f_1^{'}}^{\sim} &
        A_2 \ar@{>->}[r]^{a_{2}} \ar@{>->}[d]_{f_2^{'}}^{\sim} &
	... \\
        B_0^{'} \ar@{>->}[r]^{b^{'}_{0}} \ar@<1ex>[d]^{r_0} &
        B_1^{'} \ar@{>->}[r]^{b^{'}_{1}} \ar@<1ex>[d]^{r_1} &
        B_2^{'} \ar@{>->}[r]^{b^{'}_{2}} \ar@<1ex>[d]^{r_2} &
	... \\
        B_0 \ar@{>->}[r]^{b_{0}} \ar@{>->}@<1ex>[u]^{s_0}_{\sim} &
        B_1 \ar@{>->}[r]^{b_{1}} \ar@{>->}@<1ex>[u]^{s_1}_{\sim} &
        B_2 \ar@{>->}[r]^{b_{2}} \ar@{>->}@<1ex>[u]^{s_2}_{\sim} &
	...
	}$
\end{center}
commutative, such that additionally all maps
\begin{center}
$B^{'}_{n-1} \Sum_{A_{n-1}} A_n \ra B^{'}_{n}$ and $B^{'}_{n-1} \Sum_{B_{n-1}} B_n \ra B^{'}_n$
\end{center}
are trivial cofibrations. If we show that $\colim f^{'}_n$ and $\colim s_n$ are weak equivalences, it will follow that $\colim f_n$ is a weak equivalence. 

It suffices therefore to prove that $\colim f_n$ is a weak equivalence under the additional assumption that $f_0$ and all $B_{n-1} \Sum_{A_{n-1}} A_n \ra B_{n}$ are trivial cofibrations. The objects denoted $A = \colim A_n$ and $B = \colim B_n$ are cofibrant by CF6 (1). 

\begin{center}
$\xymatrix{
    A_0 \ar@{>->}[r]^{a_{0}} \ar@{>->}[d]_{f_0} & A_1 \ar@{>->}[r]^{a_{1}} \ar@{>->}[d] & ... \ar@{>->}[r] & A \ar@{>->}[d] \\
    B_0 \ar@{>->}[rd]_{b_{0}} \ar@{>->}[r] & B_0 \Sum_{A_0} A_1 \ar@{>->}[r] \ar@{>->}[d] & ... \ar@{>->}[r] & B_0 \Sum_{A_0} A \ar@{>->}[d] \\
    & B_1 \ar@{>->}[r] \ar@{>->}[rd]_{b_1} & ... \ar@{>->}[r] & B_1 \Sum_{A_1} A \ar@{>->}[d] \\
    & & ... \ar@{>->}[rd] & ... \ar@{>->}[d] \\
    &&& B
    }$
\end{center}
The map $A \ra B$ factors as the composition of the direct sequence of maps $A \ra B_0 \Sum_{A_0} A$ followed by $B_{n-1} \Sum_{A_{n-1}} A \ra B_{n} \Sum_{A_{n}} A$ for $n \ge 1$, and each map in the sequence is a trivial cofibration as the pushout of the trivial cofibrations $f_0$ resp. $B_{n-1} \Sum_{A_{n-1}} A_{n} \ra B_{n}$, so by CF6 (2) the map $A \ra B$ is a trivial cofibration.

The proof of statement (2) is dual to the proof of (1). 
\end{proof}

\begin{lem}[Equivalent formulation of CF6]
\label{lem:equivformaxiomcf6}
\mbox{}
\begin{enumerate}
\item
Suppose that $({\cat M}$, ${\cat W}$, ${\cat Cof})$ satisfies axioms CF1-CF4 and CF6 (1). Then the following are equivalent:
\begin{enumerate}
\item
It satisfies axiom CF6 (2)
\item
It satisfies the conclusion of \lemmaref{lem:transcompequiv} (1)
\item
It satisfies the conclusion of \lemmaref{lem:transcompequiv} (1) for any map $f_n$ of countable direct sequences of cofibrations, with each $f_n$ a trivial cofibration.
\end{enumerate}
\item
Suppose that $({\cat M}$, ${\cat W}$, ${\cat Fib})$ satisfies axioms F1-F4 and F6 (1). Then the following are equivalent:
\begin{enumerate}
\item
It satisfies axiom F6 (2)
\item
It satisfies the conclusion of \lemmaref{lem:transcompequiv} (2)
\item
It satisfies the conclusion of \lemmaref{lem:transcompequiv} (2) for any map $f_n$ of countable inverse sequences of fibrations, with each $f_n$ a trivial fibration.
\end{enumerate}
\end{enumerate}
\end{lem}

\begin{proof}
We only prove (1). The implication (a) $\Ra$ (b) is proved by \lemmaref{lem:transcompequiv}, and (b) $\Ra$ (c) is trivial. 

Let us prove (c) $\Ra$ (a). Suppose we have a countable direct sequence of trivial cofibrations with $A_0$ cofibrant
\begin{center}
$ \xymatrix {
        A_0 \ar@{>->}[r]^{a_{0}}_\sim &
        A_1 \ar@{>->}[r]^{a_{1}}_\sim &
        A_2 \ar@{>->}[r]^{a_{2}}_\sim &
	... \\
	}$
\end{center}
Then $\colim A_n$ exists and $A_0 \ra \colim A_n$ is a cofibration by axiom CF6 (1). If we view $A_0$ as a constant, countable direct sequence of identity maps, we get a map $f_n = a_{n-1} ... a_0$ of countable direct sequences of cofibrations, with each $f_n$ a trivial cofibration. From the conclusion of \lemmaref{lem:transcompequiv} (1) we see that $A_0 \ra \colim A_n$ is a weak equivalence.
\end{proof}

\section{Over and under categories}
\label{sec:catoverunder}
\index{category!over, under category}
If $u \colon {\cat A} \ra {\cat B}$ is a functor and $B$ is an object of ${\cat B}$, the over category $(u \downarrow B)$ by definition has:
\begin{enumerate}
\item
as objects, pairs $(A, g)$ of an object $A \eps {\cat A}$ and a map $g \colon uA \ra B$
\item
as maps $(A_1, g_1) \ra (A_2, g_2)$, the maps $f \colon A_1 \ra A_2$ such that $g_2 \circ uf = g_1$.
\end{enumerate}
The under category $(B \downarrow u)$ by definition has:
\begin{enumerate}
\item
as objects, pairs $(A, g)$ with $A \eps {\cat A}$ and $g \colon b \ra uA$
\item
as maps $(A_1, g_1) \ra (A_2, g_2)$, the maps $f \colon A_1 \ra A_2$ such that $uf \circ g_1 = g_2$.
\end{enumerate}

The two definitions are dual in the sense that $(B \downarrow u) \cong (u^{op} \downarrow B)^{op}$.

We have a canonical functor $i_{u, B} \colon (u \downarrow B) \ra {\cat A}$ that sends an object $(A, g: ua \ra B)$ to $A \eps {\cat A}$ and a map $(A_1, g_1) \ra (A_2, g_2)$ to the component map $A_1 \ra A_2$. Dually, we have a canonical functor $i_{B, u} \colon (B \downarrow u) \ra {\cat A}$ defined by $i_{B, u} = (i_{u^{op}, B})^{op}$.

If $A \eps {\cat A}$ is an object, for simplicity we denote $({\cat A} \downarrow A)$ for $(1_{\cat A} \downarrow A)$ and $(A \downarrow {\cat A})$ for $(A \downarrow 1_{\cat A})$. 

Suppose that ${\cat M}$ is a category and ${\cat F}$ is a class of maps of ${\cat M}$. For example, ${\cat M}$ could be a cofibration category, and ${\cat F}$ could be ${\cat W}$ or ${\cat Cof}$. For an object $A \eps {\cat M}$, we say that $i_{1_{\cat M}, A}^{-1} {\cat F}$ is the class of maps of $({\cat M} \downarrow A)$ {\it induced} by ${\cat F}$. Dually, $i_{A, 1_{\cat M}}^{-1} {\cat F}$ is the class of maps induced by ${\cat F}$ on $(A \downarrow {\cat M})$.

If $({\cat M}, {\cat W}, {\cat Cof})$ is a precofibration category, we therefore obtain an induced class of weak equivalences and cofibrations on $({\cat M} \downarrow A)$ and on $(A \downarrow {\cat M})$. Dually if $({\cat M}, {\cat W}, {\cat Fib})$ is a prefibration category we obtain an induced class of weak equivalences and fibrations on $({\cat M} \downarrow A)$ and on $(A \downarrow {\cat M})$.

The following result is a simple consequence of the definitions.

\begin{prop}
\label{prop:overcofcat}
\mbox{}
\begin{enumerate}
\item
Suppose that ${\cat M}$ is a (pre)cofibration category, and $A \eps {\cat M}$ is an object. Then:
\begin{enumerate}
\item
$({\cat M} \downarrow A)$ is a (pre)cofibration category
\item
If $A$ is cofibrant, $(A \downarrow {\cat M})$ is a (pre)cofibration category
\end{enumerate}
with respect to the induced weak equivalences and cofibrations.
\item
Suppose that ${\cat M}$ is a (pre)fibration category, and $A \eps {\cat M}$ is an object. Then:
\begin{enumerate}
\item
$(A \downarrow {\cat M})$ is a (pre)fibration category
\item
If $A$ is fibrant, $({\cat M} \downarrow A)$ is a (pre)fibration category 
\end{enumerate}
with respect to the induced weak equivalences and fibrations. $\square$
\end{enumerate}
\end{prop}

\section{Properness}
\label{sec:properness}
Sometimes, a precofibration category $({\cat M}, {\cat W}, {\cat Cof})$ admits more than one precofibration structure with weak equivalences ${\cat W}$. Under an additional condition (left properness) one can show that $({\cat M}, {\cat W})$ admits an intrinsic structure of a CF1-CF4 precofibration category, larger than ${\cat Cof}$, defined in terms of what we will call {\it left proper} maps. 

\begin{defn}
\label{defn:propercofcat}
\mbox{}
\begin{enumerate}
\item
A precofibration category $({\cat M}, {\cat W}, {\cat Cof})$ is {\it left proper} if it satisfies
\begin{description}
\item[PCF] 
Given a solid diagram in ${\cat M}$ with $i$ a cofibration and $A$ cofibrant,
\begin{center}
$ \xymatrix {
	A \ar[r]^r \ar@{>->}[d]_i &
	C \ar@{>-->}[d] \\
	B \ar@{-->}[r]^{r^{'}} &
	D
	}$
\end{center}
then the pushout exists in ${\cat M}$. Moreover, if $r$ is a weak equivalence then so is $r^{'}$.
\end{description}
\item
A prefibration category $({\cat M}, {\cat W}, {\cat Fib})$ is {\it right proper} if 
\begin{description}
\item[PF] 
Given a solid diagram in ${\cat M}$ with $p$ a fibration and $A$ fibrant,
\begin{center}
$ \xymatrix {
        D \ar@{-->}[r]^{r^{'}} \ar@{-->}[d] & B \ar@{->>}[d]^p \\
        C \ar[r]^r_{\sim} & A
	}$
\end{center}
then the pullback exists in ${\cat M}$. Moreover, if $r$ is a weak equivalence then so is $r^{'}$.
\end{description}
\item
An ABC model category is proper if its underlying precofibration and prefibration categories are left resp. right proper.
\end{enumerate}
\end{defn}

From the Excision Lemma, a precofibration category with all objects cofibrant is left proper. A prefibration category with all objects fibrant is right proper.

We will define proper maps in the context of what we call {\it category pairs}.

\index{category!with weak equivalences}
\begin{defn}
\label{defn:catweakequiv}
A category pair $({\cat M}$, ${\cat W})$ consists of a category ${\cat M}$ with a class of weak equivalence maps ${\cat W}$, where ${\cat W}$  is stable under composition and includes the identity maps of ${\cat M}$. We may view ${\cat W}$ as defining a subcategory with the same objects as ${\cat M}$. 
\end{defn}

\begin{defn}
\label{defn:propercofibration}
Suppose that $({\cat M}, {\cat W})$ is a category pair.
\begin{enumerate}
\item
A map $f \colon A \ra B$ is called {\it left proper} if for any diagram of full maps with $r$ a weak equivalence
\begin{center}
$\xymatrix {
    A \ar[r] \ar[d]_f & C_1 \ar[r]^r_{\sim} \ar@{-->}[d] & C_2 \ar@{-->}[d] \\
    B \ar@{-->}[r] & D_1 \ar@{-->}[r]^{r^{'}} & D_2
  }$
\end{center}
the pushouts exist, and the map $r^{'}$ is again a weak equivalence. 
\item
A map $f \colon B \ra A$ is called {\it right proper} if for any diagram of full maps with $r$ a weak equivalence
\begin{center}
$\xymatrix {
    D_2 \ar@{-->}[r]^{r^{'}} \ar@{-->}[d] & D_1 \ar@{-->}[r] \ar@{-->}[d] & B \ar[d]^f \\
    C_2 \ar[r]^r_{\sim} & C_1 \ar[r] & A
  }$
\end{center}
the pullbacks exist, and the map $r^{'}$ is again a weak equivalence. 
\end{enumerate}
\end{defn}

We say that an object $A$ is {\it left proper} if the map $\initial \ra A$ is left proper. An object $A$ is {\it right proper} if the map $A \ra \terminal$ is right proper.

The class of left proper maps of $({\cat M}, {\cat W})$ will be denoted ${\cat PrCof}$. The class of right proper maps will be denoted ${\cat PrFib}$. Observe that the left proper maps are stable under composition and under pushout. The right proper maps are stable under composition and under pullback. All isomorphisms are left and right proper.

The left proper weak equivalences will be called {\it trivial left proper} maps, and the right proper weak equivalences will be called {\it trivial right proper} maps.

\begin{thm}
\label{thm:leftproperstruct}
\mbox{}
\begin{enumerate}
\item
If $({\cat M}, {\cat W}, {\cat Cof})$ is a left proper precofibration category, then 
\begin{enumerate}
\item
Any cofibration $A \ra B$ with $A$ cofibrant is left proper.
\item
Any map of ${\cat M}$ factors as a left proper map followed by a weak equivalence.
\item
Trivial left proper maps are stable under pushout.
\item
$({\cat M}, {\cat W}, {\cat PrCof})$ is a precofibration category.
\end{enumerate}
\item
If $({\cat M}, {\cat W}, {\cat Fib})$ is a right proper prefibration category, then 
\begin{enumerate}
\item
Any fibration $A \ra B$ with $B$ fibrant is right proper.
\item
Any map of ${\cat M}$ factors as a weak equivalence followed by a right proper map.
\item
Trivial right proper maps are stable under pushout.
\item
$({\cat M}, {\cat W}, {\cat PrFib})$ is a prefibration category.
\end{enumerate}
\end{enumerate}
\end{thm}

\begin{proof}
We only prove (1). For (1) (a), suppose we have a diagram
\begin{center}
$\xymatrix {
    A \ar[r]^f \ar@{>->}[d]_i & C_1 \ar[r]^r_{\sim} \ar[d] & C_2 \ar[d] \\
    B \ar[r] & D_1 \ar[r]^{r^{'}} & D_2
  }$
\end{center}
with $A$ cofibrant, $i$ a cofibration and $r$ a weak equivalence, such that both squares are pushouts. We'd like to show that $r^{'}$ is a weak equivalence. 

Denote $f = f^{'}s$ a factorization of $f$ as a cofibration $f^{'}$ followed by a weak equivalence $s$. 
\begin{center}
$\xymatrix {
    A \ar@{>->}[r]^{f_1} \ar@{>->}[d]_i & C^{'}_1 \ar[r]^s_\sim \ar@{>->}[d]^{i^{'}} & C_1 \ar[r]^r_{\sim} \ar[d] & C_2 \ar[d] \\
    B \ar@{>->}[r] & D^{'}_1 \ar[r]^{s^{'}} & D_1 \ar[r]^{r^{'}} & D_2
  }$
\end{center}
Denote $i^{'}$ the pushout of $i$, and $s^{'}$ the pushout of $s$. Since $i^{'}$ is a cofibration and $s$, $rs$ are weak equivalences, from the left properness of ${\cat M}$ we see that $s^{'}$, $r^{'}s^{'}$ and therefore $r^{'}$ are weak equivalences.

For (1) (b), suppose $f \colon A \ra B$ is a map in ${\cat M}$. We'd like to construct a factorization $f = rf^{'}$ with $f^{'}$ left proper and $r$ a weak equivalence.

Let $a \colon A^{'} \ra A$ be a cofibrant replacement of $A$, and $fa = r^{'}_1f_1$ factorization of $fa$ as a cofibration $f_1$ followed by a weak equivalence $r_1$.
\begin{center}
$\xymatrix{ 
    &&& B \\
    A \ar[rr]_{f^{'}} \ar[urrr]^f && Y^{'} \ar[ur]^(.3)r & \\
    A^{'} \ar[u]_\sim^{a^{'}} \ar@{>->}[rr]_{f_1} && A_1 \ar[u]^{b^{'}}_\sim \ar[uur]_{r_1^{'}}^\sim &
  }$
\end{center}
Define $f^{'}$ as the pushout of $f_1$. The map $b^{'}$ is a weak equivalence as the pushout of $a^{'}$, since ${\cat M}$ is left proper. The map $r$ is a weak equivalence by the two out of three axiom, and the map $f^{'}$ is left proper as the pushout of $f_1$ which is a cofibration with cofibrant domain (therefore left proper).

For (1) (c), let $i \colon A \ra B$ be a trivial left proper map and let $f \colon A \ra C$ be a map. In the diagram
\begin{center}
$\xymatrix {
    A \ar[r]^{f^{'}} \ar[d]_i^\sim & C^{'} \ar[r]^r_{\sim} \ar[d]^{i^{'}} & C \ar[d]^{j} \\
    B \ar[r] & D^{'} \ar[r]^{r^{'}} & D
  }$
\end{center}
using part (b) we have factored $f$ as $rf^{'}$, with $f^{'}$ left proper and $r$ a weak equivalence. We define $i^{'}$ and $j$ as the pushouts of $i$. We therefore have that $i^{'}$ and $j$ are left proper. The map $i^{'}$ is a weak equivalence since $f^{'}$ is left proper. The map $r^{'}$ is a weak equivalence since $i$ is proper. From the two out of three axiom, the map $j$ is also a weak equivalence, therefore a trivial left proper map.

For (1) (d), the axioms CF1, CF2 and CF3 (1) are trivially verified. Part (c) proves axiom CF3 (2), and part (b) proves axiom CF4.
\end{proof}

It does not appear to be the case that $({\cat M}, {\cat W}, {\cat PrCof})$ necessarily satisfies CF5 or CF6 if $({\cat M}, {\cat W}, {\cat Cof})$ does.

In the rest of the section, we will review briefly a number of elementary properties of proper maps.

\begin{prop}
\label{prop:propermaps1}
\mbox{}
\begin{enumerate}
\item
Suppose that $({\cat M}, {\cat W}, {\cat Cof})$ is a left proper precofibration category. Then a weak equivalence is a trivial left proper map iff all its pushouts exist and are weak equivalences.
\item
Suppose that $({\cat M}, {\cat W}, {\cat Fib})$ is a right proper prefibration category. Then a weak equivalence is a trivial right proper map iff all its pullbacks exist and are weak equivalences.
\end{enumerate}
\end{prop}

\begin{proof}
We only prove (1). Implication $\Ra$ follows from \theoremref{thm:leftproperstruct} (1) (c). 

For $\La$, suppose that $i \colon A \ra B$ is a weak equivalence whose pushouts remain weak equivalences. We'd like to show that $i$ is left proper. For any map $f$ and weak equivalence $r$, we construct the diagram with pushout squares
\begin{center}
$\xymatrix {
    A \ar[r]^f \ar[d]_i^\sim & C_1 \ar[r]^r_{\sim} \ar[d]^{i_1}_\sim & C_2 \ar[d]^{i_2}_\sim \\
    B \ar[r] & D_1 \ar[r]^{r^{'}} & D_2
  }$
\end{center}
The maps $i_1$, $i_2$ and are weak equivalences as pushouts of $i$. From the 2 out of 3 axiom, the map $r^{'}$ is a weak equivalence, which shows that $i$ is left proper.
\end{proof}

\begin{prop}
\label{prop:propermaps2}
\mbox{}
\begin{enumerate}
\item
Suppose that $({\cat M}, {\cat W}, {\cat Cof})$ is a cocomplete, left proper precofibration category. If $f, g$ are two composable maps such that $gf$ is left proper and $g$ is trivial left proper. Then $f$ is left proper.
\item
Suppose that $({\cat M}, {\cat W}, {\cat Fib})$ is a complete, right proper prefibration category. If $f, g$ are two composable maps such that $gf$ is left proper and $g$ is trivial left proper. Then $f$ is left proper.
\end{enumerate}
\end{prop}

\begin{proof}
We only prove (1). In the diagram
\begin{center}
$\xymatrix {
    A_1 \ar[r] \ar[d]_f & B_1 \ar[r]^r_{\sim} \ar[d] & C_1 \ar[d] \\
    A_2 \ar[r] \ar[d]_g^\sim & B_2 \ar[r]^{r^{'}} \ar[d]^{g^{'}} & C_2 \ar[d]^{g^{''}} \\
    A_3 \ar[r] & B_3 \ar[r]^{r^{''}} & C_3
  }$
\end{center}
$r$ is a weak equivalence and all squares are pushouts. To prove that $f$ is left proper, we need to show that $r^{'}$ is a weak equivalence. But $r^{''}$ is a weak equivalence since $gf$ is proper. $g^{'}$ and $g^{''}$ are trivial left proper as pushouts of $g$. It follows from the two out of three axiom that $r^{'}$ is a weak equivalence. 
\end{proof}

Recall the definition of the retract of a map
\begin{defn}
\label{defn:mapretract}
\index{retract}
A map $f \colon A \ra B$ in a category ${\cat M}$ is a retract of $g \colon C \ra D$ if there exists a commutative diagram
\begin{center}
$\xymatrix{
	A \ar[r] \ar[d]_f \ar@/^1pc/[rr]^{1_A} &
	C \ar[d]^g \ar[r] & A \ar[d]^f \\
	B \ar[r] \ar@/_1pc/[rr]_{1_B} &
	D \ar[r] & B
	}$
\end{center}
\end{defn}

Note that the saturation $\overline{\cat W}$ of the class of weak equivalences ${\cat W}$ is closed under retracts.

\begin{prop}
\label{prop:propermaps3}
Suppose that $({\cat M}, {\cat W})$ is a category pair and that ${\cat W}$ is closed under retracts.
\begin{enumerate}
\item
Assume that ${\cat M}$ is cocomplete. Then the class of left proper maps and that of trivial left proper maps are both closed under retracts.
\item
Assume that ${\cat M}$ is complete. Then the class of right proper maps and that of trivial right proper maps are both closed under retracts.
\end{enumerate}
\end{prop}

\begin{proof}
Follows directly from the definitions.
\end{proof}
